\chapter{Quality control and acceptance sampling}
\label{ch:calidad}

\section{Monitoring a process or accepting a lot}
Statistical control looks for signals of change in a process; acceptance
sampling makes decisions about lots using samples. A stable process can produce
defects, and accepting a lot does not certify that all its units conform.
Both procedures require consistent definitions and sampling
\cite{montgomery2020quality}.

QSB supports c, p, and \(\bar X\)-R charts, Pareto charts, and normal
probability plots; it also evaluates single and double binomial acceptance
plans. The native backend is deterministic, and the validation backend does
not solve. No external backend is available for these families.

\section{c and p charts}
The c chart counts defects in an inspection unit of constant size. An item
may have more than one defect. Under a Poisson count assumption,
\[
 CL=\bar c,\qquad LCL=\max(0,\bar c-3\sqrt{\bar c}),\quad
 UCL=\bar c+3\sqrt{\bar c}.
\]
\path{manual/examples/prob_calidad.json} contains counts
\((2,3,4,3,2,4,3,3,2,4)\). From the root:
\begin{lstlisting}
.build/debug/qsb validate-quality-json manual/examples/prob_calidad.json
.build/debug/qsb solve-quality-json manual/examples/prob_calidad.json
\end{lstlisting}
The center line is \(3\), the lower limit zero, and the upper limit
\(8.196152\). There are no points outside the limits or signals from the
implemented rules. This is an educational check of the formula; ten counts
are insufficient to establish an industrial process as stable. Inspection
units of varying sizes require another approach; QSB does not implement a u
chart in this route.

The p chart uses the proportion of defective units, with size \(n_i\) for
each subgroup. Its weighted center is
\[
 \bar p=\frac{\sum_i n_ip_i}{\sum_i n_i},\qquad
 CL_i=\bar p,\quad CL_i\pm3\sqrt{\bar p(1-\bar p)/n_i},
\]
with limits clipped to \([0,1]\).

In \path{prob_p.json}, sizes
\((100,100,200)\) and proportions \((0.02,0.04,0.03)\) produce
\(\bar p=0.03\); upper limits are \(0.081176\) for the first two
subgroups and \(0.066187\) for the third. The lower limit is zero.
Input uses fractions: \(0.03\) means \(3\%\), rather than \(0.03\%\).
Three-sigma limits are a statistical approximation to the data distribution,
even though QSB labels the application of the formulas as \texttt{exact}.

\section{Mean and range charts}
Rational subgroups combine measurements taken under comparable conditions.
With \(\bar{\bar X}\) as the mean of subgroup means and \(\bar R\)
as the mean of ranges,
\[
 LCL_{\bar X},UCL_{\bar X}=\bar{\bar X}\mp A_2\bar R,\qquad
 LCL_R=D_3\bar R,\quad UCL_R=D_4\bar R.
\]
QSB uses standard constants for uniform sizes between two and ten. Review
variation in R before interpreting shifts in means.
\path{prob_xr.json} contains subgroups \((9,10,11)\), \((10,11,12)\),
and \((8,9,10)\). The results are \(\bar{\bar X}=10\), \(\bar R=2\),
mean limits \([7.954,12.046]\), and range limits \([0,5.148]\), with no
points outside the limits. Control limits describe process variability; they
are not design tolerances. QSB does not calculate Cp/Cpk indices in this
workflow.

\section{Pareto and normal probability}
Pareto aggregates categories, sorts them by descending frequency, and reports
shares and cumulative shares. \path{prob_pareto.json} has rows \((4,2,1)\)
and \((2,1,0)\) for Rayas (scratches), Manchas (stains), and Roturas
(breakages). Their totals are \((6,3,1)\), the overall total is \(10\),
shares are \((0.6,0.3,0.1)\), and cumulative shares are \((0.6,0.9,1)\).
Fields named \texttt{percentage} in JSON store fractions, rather than values
from zero to one hundred. Prioritizing Rayas by frequency does not establish
that it is the category with the greatest cost or risk.

The normal plot sorts measurements and assigns Blom plotting positions
\(p_i=(i-3/8)/(n+1/4)\); it uses \(z_i=\Phi^{-1}(p_i)\), fitting
measurement against normal score. \path{prob_normal.json} contains
\((8,9,10,11,12)\); mean \(10\), sample standard deviation \(1.581139\),
intercept \(10\), slope \(1.742919\), and correlation \(0.997902\).
It is labeled \texttt{approximate} because of the approximate normal quantile.
High correlation with five points does not demonstrate normality; this is not
a test with a p-value. These four examples run using the same quality commands
with the file changed.

\section{Signals and the quality interface}
QSB flags points beyond three sigma and fixed Western Electric rules: two of
three beyond two sigma on the same side, four of five beyond one sigma on the
same side, and eight on the same side of the center. Points already flagged
beyond three sigma are excluded from the counts for the two-of-three and
four-of-five zone rules. There is no rule configuration or normalized metadata
for cause, action, and comment. Signal indices identify observations to
investigate, rather than automatic cause diagnoses. Clipped attribute-chart
limits affect the scale used by this implementation's rules: retain the output
and review each signal alongside the chart.

In macOS, select the variant in the native editor and complete its table:
counts/proportions, subgroups, categories, or measurements. Changing variants
preserves inactive tables during the session, but JSON exports only the active
one. Paste complete tables to replace them atomically; incomplete imported
data remain as invalid cells. After validation and solving, charts show centers,
limits, and signals; Pareto shows frequency and cumulative share; the normal
plot shows points and a fitted line. Consult tables and JSON for the precise
values of the calculation, and preserve acquisition order.

\section{Single acceptance plan}
For sample size \(n\) and maximum acceptable number of defective units
\(c\), the lot is accepted if \(X\le c\). Under independent binomial
sampling with proportion \(p\), the operating characteristic curve is
\[
 P_a(p)=\sum_{x=0}^c{n\choose x}p^x(1-p)^{n-x}.
\]
AQL is the acceptable quality level and RQL the rejectable quality level.
Actual producer's risk is \(\alpha=1-P_a(AQL)\); consumer's risk is
\(\beta=P_a(RQL)\). Entering nominal risks does not alter an already
supplied plan to satisfy them.

\path{prob_muestreo.json} sets \(n=20,c=1,AQL=0.01,RQL=0.15\), nominal
risks \(0.05,0.10\), and lot size \(500\).
\begin{lstlisting}
.build/debug/qsb validate-acceptance-json manual/examples/prob_muestreo.json
.build/debug/qsb solve-acceptance-json manual/examples/prob_muestreo.json
\end{lstlisting}
QSB returns \(P_a(AQL)=0.983141\), producer's risk \(0.016859\), and
consumer's risk \(0.175558\). The plan exceeds the nominal consumer's risk:
it must not be described as satisfying \(\beta=0.10\). Output contains
101 curve points from \(0\) to \(100\%\) defective, as well as measures
at AQL and RQL. This grid does not replace continuous evaluation when optimizing
a risk between points.

Average sample number ASN is \(20\). Under rectifying inspection of rejected
lots and removal of inspected defective units,
\[
 ATI=n+(N-n)(1-P_a),\qquad AOQ=pP_a(N-n)/N.
\]
At AQL, ATI is \(28.092482\) and AOQ \(0.009438\), that is
\(0.943815\%\). If rejected lots are returned without full inspection,
these measures do not describe the actual procedure. QSB evaluates binomial
probabilities exactly within the assumption, but sampling a large fraction
of the lot without replacement requires considering the hypergeometric
distribution; the core provides that probability, rather than a selector in
this command.

\section{Double plans and design}
In a double plan, \(n_1\) units are inspected. If the number of defective
units is \(x_1\le c_1\), the lot is accepted; if \(x_1\ge r_1\), it is
rejected; in between, \(n_2\) \emph{additional} units are inspected. Final
acceptance occurs if \(x_1+x_2\le c_2\). Confusing \(n_2\) with cumulative
sample size changes the plan. Its acceptance probability is
\[
 P_a=P(X_1\le c_1)+
 \sum_{x=c_1+1}^{r_1-1}P(X_1=x)P(X_2\le c_2-x).
\]
The file \path{prob_doble.json} uses \(n_1=10,c_1=0,r_1=3\),
\(n_2=10,c_2=1\), with the same quality levels and lot size as before.
Changing the file in \texttt{solve-acceptance-json} gives producer's risk
\(0.013001\), consumer's risk \(0.265274\), ASN at AQL \(10.955041\),
and ASN at RQL \(16.233221\). It reduces average inspection at the cost
of accepting more bad lots in this example; it is not a universal improvement.
ATI at AQL is \(17.313145\) and AOQ \(0.009654\), under the solver's
rectifying convention. For double plans, QSB replaces \(n\) with ASN in
the reported ATI and AOQ formulas. This is an aggregate convention: it does
not separately calculate the inspected size conditional on each acceptance
or rejection route.

The CLI can design the smallest \emph{single} plan satisfying binomial bounds:
\begin{lstlisting}
.build/debug/qsb design-acceptance 500 0.01 0.15 0.05 0.10
\end{lstlisting}
The result is \(n=25,c=1\), producer's risk \(0.025759\), and consumer's
risk \(0.093071\). Argument order is lot size, AQL, RQL, producer's risk,
and consumer's risk. The command searches for a plan, rather than evaluating
one already fixed. It does not design double plans, optimize expected cost,
implement nonzero classification errors, or claim conformity with a specific
sampling standard.

In macOS, the form distinguishes single/double plans and their decision
thresholds; it preserves shared quality levels, risks, and economic data when
the type changes. Validate and solve to see the OC curve, AQL/RQL markers,
actual versus nominal risks, and ASN/ATI/AOQ. Economic fields are preserved
in the model; core output contains quality and inspection measures, rather
than an economic optimum. JSON allows the exact plan used to be audited.

\section{Exercises}
\begin{enumerate}
\item Add a count of twenty to the c example and observe how the center and
limits change when it is included in estimation. Discuss when to establish a
reference phase.
\item Compare p charts with sizes of one hundred and one thousand, keeping
proportions fixed. Explain why the same deviation can produce different signals.
\item Design a single plan using the example's bounds and compare ASN with
the \(20,1\) plan. Check actual risks, rather than only nominal ones.
\end{enumerate}
